\documentclass{article}
% Source: Topic_06_Teacher_Edition.pdf, PDF pages 99-100 (printed pages 341A-341B): Topic Assessment Form A with the TE's marked answers, Form B with marked answers, Item Analysis for both forms. Cross-checked against the Savvas Realize Form B PDF text layer (Topic6AssessAnswer.pdf), the Savvas Form B docx (Topic6Assess.docx, math dropped) and the earlier Lesson_planning/Topic6AssessFormB.tex.
\usepackage{amsmath}
\usepackage{amssymb}
\newenvironment{assessment-meta}{}{}   % topic, title, forms, item count, printed pages, sources
\newenvironment{item-analysis}{}{}     % Savvas's Item / DOK / Skills Review and Practice table (same for both forms)
\newenvironment{ta-item}[2]{}{}        % #1 = item number, #2 = form (A or B)
\newcommand{\answer}[1]{}              % the TE's marked answer, inside a ta-item
\newcommand{\placeholder}[2]{}         % #1 = type (graph, table, figure, number-line), #2 = complete description
\newcommand{\uncertain}[2]{#1}         % #1 = best reading, #2 = alternatives and why
\newcommand{\te}[1]{}                  % transcriber's note (a discrepancy, a publisher misprint, a computed check)
\begin{document}
\begin{assessment-meta}
topic: 6
title: Exponential and Logarithmic Functions
forms: A and B (parallel item for item)
items: 17
printed: 341A-341B
sources: Topic_06_Teacher_Edition.pdf, PDF page 99 (341A, Form A, both student pages, answers marked in magenta) and PDF page 100 (341B, Form B, both student pages, answers marked; Item Analysis); cross-checked against Savvas Realize Topic6AssessAnswer.pdf (Form B text layer), Topic6Assess.docx (Form B, math dropped) and Lesson_planning/Topic6AssessFormB.tex
\end{assessment-meta}
\te{Student page headers read "6 Topic Assessment Form A" and "6 Topic Assessment Form B"; page footers "enVision Algebra 2 | 1of2 / 2of2 | Assessment Sourcebook". TE running head "TOPIC 6 Exponential and Logarithmic Functions", banner "TOPIC ASSESSMENT Form A / Form B", footers "TOPIC 6 | 341A | TOPIC ASSESSMENT" and "TOPIC 6 | 341B | TOPIC ASSESSMENT". No running-head misprint on these pages.}
\te{TE note under Form A (341A): "Topic Assessment Assess students' understanding of topic concepts and skills using the Topic Assessment found at SavvasRealize.com. These auto-scored online assessments provide students with a breadth of technology-enhanced item types. There are two versions of the Topic Assessment, Form A and Form B. These two versions, available in print and at SavvasRealize.com, are parallel tests that assess the same content item for item. The Item Analysis chart on the next page can be used for both versions."}
% Source page 100: 341B Item Analysis
\begin{item-analysis}
Item & DOK & Skills Review and Practice \\
1 & 2 & F11 \\
2 & 2 & F14 \\
3 & 3 & A08 \\
4 & 2 & F13 \\
5 & 2 & F13, F14 \\
6 & 2 & F14 \\
7 & 2 & F14 \\
8 & 2 & F69 \\
9 & 2 & F82 \\
10 & 2 & F71 \\
11 & 2 & F69, F70 \\
12 & 2 & F70 \\
13 & 2 & F69 \\
14 & 2 & F71 \\
15 & 2 & F72 \\
16 & 2 & F82 \\
17 & 3 & F71 \\
\end{item-analysis}
% Source page 99: 341A Topic Assessment Form A
\begin{ta-item}{1}{A}
Which function (g) represents the exponential function (f(x)=5^x) after a vertical stretch by a factor of 2 and a reflection across the (x)-axis?
A. (g(x)=(2\cdot 5)^{-x})
B. (g(x)=-(2\cdot 5)^x)
C. (g(x)=2\cdot(5)^{-x})
D. (g(x)=-2\cdot(5)^x)
\answer{1A. D}
\end{ta-item}
\begin{ta-item}{2}{A}
A population of butterflies increases from 2000 to 2800 over 7 years. Which function best models the population (x) years from the first measurement?
A. (f(x)=2000(1+0.4)^{\frac{1}{7}x})
B. (f(x)=2000(1+0.4^7)^x)
C. (f(x)=2000\left(1+\frac{0.4}{7}\right)^x)
D. (f(x)=2000\left(1+\frac{0.4}{7}\right)^{7x})
\answer{2A. A}
\end{ta-item}
\begin{ta-item}{3}{A}
In the expression shown (t\neq s).
(\frac{(s-t)+(2s-2t)^2}{t-s})
When (x) replaces ((s-t)) which expression is equivalent to the expression shown. Select all that apply.
A. (\frac{x+4x^2}{-x})
B. (\frac{x-2x^2}{x})
C. (\frac{x+(4x)^2}{x})
D. (\frac{-x-4x^2}{x})
E. (\frac{x+(2x)^2}{x})
\answer{3A. A, D}
\end{ta-item}
\begin{ta-item}{4}{A}
Which of the following functions has a faster decay factor than the function shown in the graph?
\placeholder{graph}{Prompt graph (not part of the answer): a decreasing exponential curve on a gray grid. Axes labeled x and y with arrowheads, origin O. x-axis labels 2 and 4, gridlines every 1 unit, window about x = -2 to 6. y-axis labels 20 and 40, gridlines every 10 units, window about y = -20 to 60. Three labeled points on the curve: (0, 48), (1, 36), (2, 27). The curve ends at the upper left in an arrowhead just inside the top of the grid, near (-0.5, 56), next to the y-axis arrow, and falls toward the x-axis to the right, ending in an arrowhead near (5.5, 10); it stays above the x-axis.}
A. (f(x)=48(0.6)^x)
B. (f(x)=48(0.75)^x)
C. (f(x)=48(0.772)^x)
D. (f(x)=48(0.8)^x)
\answer{4A. A}
\end{ta-item}
\begin{ta-item}{5}{A}
The number of views (y) on a viral video after (x) days can be represented by an exponential model, (y=a\cdot b^x). After 1 day, a cat video has 10 views, and after 2 days it has 20 views. Which represents another point on the graph of this model?
A. ((3, 25))
B. ((3, 30))
C. ((3, 40))
D. ((3, 60))
\answer{5A. C}
\end{ta-item}
\begin{ta-item}{6}{A}
A warren of rabbits has a population of 25. The population is increasing at a rate of 20\% per year. Select the model of the population after (t) years that reveals the monthly growth rate.
A. (y=25(1.2)^t)
B. (y=25(1.0167)^{12t})
C. (y=25(1.0167)^t)
D. (y=25(1.0153)^{12t})
\answer{6A. D}
\end{ta-item}
\begin{ta-item}{7}{A}
What function is the inverse of the exponential function (y=4^x)?
A. (y=x^{\frac{1}{4}})
B. (y=\left(\frac{1}{4}\right)^x)
C. (y=\log_x 4)
D. (y=\log_4 x)
\answer{7A. D}
\end{ta-item}
\begin{ta-item}{8}{A}
What is the value of (n) in the formula (S=\frac{a_1(1-r^n)}{(1-r)}) for the sum of the geometric series (5+3\frac{1}{3}+2\frac{2}{9}+\ldots+\frac{640}{2187})?
(n=) [ ]
\answer{8A. 8}
\end{ta-item}
\begin{ta-item}{9}{A}
What is the solution to the equation (\log_4(5x+9)=3)?
(x=) [ ]
\answer{9A. 11}
\end{ta-item}
\begin{ta-item}{10}{A}
The graph shows the function (f(x)=6(3)^x). What is the value of the inverse function, (f^{-1}), at (x=2)?
\placeholder{graph}{Prompt graph (not part of the answer), to the right of the text: an increasing exponential curve on a gray grid. Axes labeled x and y with arrowheads, origin O. x-axis labels -4, -2, 2, gridlines every 1 unit, window about x = -6 to 4. y-axis labels 4, 8, 12, 16, gridlines every 2 units, window about y = -4 to 20. Labeled points on the curve: (-2, 2/3) and (-1, 2) (both labels with leader lines), (0, 6), (1, 18). The curve approaches the x-axis (y = 0) to the left, ending in an arrowhead, and rises steeply off the top of the grid to the right of the y-axis, ending in an arrowhead.}
(f^{-1}(2)=) [ ]
\answer{10A. -1}
\end{ta-item}
\begin{ta-item}{11}{A}
What is the asymptote of the graph of (g(x)=\ln x+5)?
A. vertical at (x=e)
B. vertical at (x=0)
C. horizontal at (y=e)
D. horizontal at (y=5)
\answer{11A. B}
\end{ta-item}
\begin{ta-item}{12}{A}
Complete the equation of the inverse of the function (f(x)=\log(2x)).
\placeholder{figure}{Tile bank of four drag tiles, left to right: 1/2, 2, 10, 1/10.}
(f^{-1}(x)=) [ ]([ ])^x
\answer{12A. (f^{-1}(x)=\frac{1}{2}(10)^x): first box 1/2, second box 10}
\end{ta-item}
\begin{ta-item}{13}{A}
Use the properties of logarithms to choose the expression (2\log 5+5\log x) as a single logarithm.
A. (\log(10+5x))
B. (\log(10x^5))
C. (\log(25x^5))
D. (\log(25+x^5))
\answer{13A. C}
\end{ta-item}
\begin{ta-item}{14}{A}
A new kind of rocket takes off with exponential acceleration with height in miles represented by (3^x), where (x) is time in minutes. Find the time when the rocket's height is 8 miles. Select the exact solution as a logarithm, and select an approximate solution.
A. (x=\log_8 3)
B. (x=\log_3 8)
C. (x=\log 8)
D. (x=1.893)
E. (x=0.528)
F. (x=0.903)
\answer{14A. B, D}
\te{Choices print in two columns: A, B, C on the left and D, E, F on the right.}
\end{ta-item}
\begin{ta-item}{15}{A}
What is the sum of the geometric series (S=\frac{27}{64}+\frac{81}{128}+\frac{243}{256}+\ldots+\frac{19{,}683}{4096})? Round to the nearest thousandth.
(S=) [ ]
\answer{15A. 13.573}
\te{computed: a_1 = 27/64, r = 3/2, n = 7, S = 55593/4096 = 13.5725..., agrees with the key.}
\end{ta-item}
\begin{ta-item}{16}{A}
Complete the steps to write the expression (\log a+3\log b-2\log c) as a single logarithm using the properties of logarithms.
\placeholder{figure}{Tile bank of three drag tiles, left to right: Product, Quotient, Power.}
(\log a+3\log b-2\log c)
(=\log a+\log b^3-\log c^2) [ ]
(=\log ab^3-\log c^2) [ ]
(=\log\left(\frac{ab^3}{c^2}\right)) [ ]
\answer{16A. Power; Product; Quotient (boxes top to bottom)}
\end{ta-item}
\begin{ta-item}{17}{A}
For the function (f(x)=-4\cdot(0.3)^x+1), identify the (y)-intercept and asymptote.
(y)-Intercept: [ ]
Asymptote: (y=) [ ]
\answer{17A. y-Intercept: -3; Asymptote: y = 1}
\end{ta-item}
% Source page 100: 341B Topic Assessment Form B
\begin{ta-item}{1}{B}
Which function (g) represents the exponential function (f(x)=3^x) after a horizontal stretch by a factor of 4 and a reflection across the (y)-axis?
A. (g(x)=(3)^{-\frac{x}{4}})
B. (g(x)=(3)^{-4x})
C. (g(x)=-(3)^{\frac{x}{4}})
D. (g(x)=-(3)^{4x})
\answer{1B. A}
\te{The Form B docx prints choice B as "g(x) = −(3)–4x" (with a leading minus); the TE page and the Realize PDF text layer have no minus: (g(x)=(3)^{-4x}). TE image followed.}
\end{ta-item}
\begin{ta-item}{2}{B}
A population of mice increases from 400 to 1400 over 7 years. Which function best models the population (x) years from the first measurement?
A. (f(x)=400\left(1+2.5^{\frac{1}{7}}\right)^x)
B. (f(x)=400(1+2.5)^{\frac{1}{7}x})
C. (f(x)=400\left(1+\frac{2.5}{7}\right)^x)
D. (f(x)=400\left(1+\frac{2.5}{7}\right)^{7x})
\answer{2B. B}
\end{ta-item}
\begin{ta-item}{3}{B}
In the expression shown \uncertain{(t\neq s)}{the not-equal glyph prints slanted and blurred on the Form B page; Form A and the Form B docx read the same span as t \neq s}.
(\frac{(s-t)+(2s-2t)^2}{t-s})
When (x) replaces ((s-t)) which expression is equivalent to the expression shown. Select all that apply.
A. (\frac{x+4x^2}{-x})
B. (\frac{x-2x^2}{x})
C. (\frac{x+(4x)^2}{x})
D. (\frac{-x-4x^2}{x})
E. (\frac{x+(2x)^2}{x})
\answer{3B. A, D}
\te{Item 3 is printed identically on both forms.}
\end{ta-item}
\begin{ta-item}{4}{B}
Which of the following functions has a slower growth rate than the function shown in the graph?
\placeholder{graph}{Prompt graph (not part of the answer): an increasing exponential curve on a gray grid. Axes labeled x and y with arrowheads, origin O. x-axis labels -4, -2, 2, 4, gridlines every 1 unit, window about x = -5 to 5. y-axis labels -20, -10, 10, 20, gridlines every 5 units, window about y = -25 to 25. Labeled points on the curve: (-1, 2) and (1, 18); the curve crosses the y-axis at about (0, 6) (unlabeled). The curve approaches the x-axis (y = 0) to the left, ending in an arrowhead, and rises off the top of the grid just right of x = 1, ending in an arrowhead.}
A. (f(x)=4(4.5)^x)
B. (f(x)=6(3.5)^x)
C. (f(x)=8(3.1)^x)
D. (f(x)=10(2)^x)
\answer{4B. D}
\end{ta-item}
\begin{ta-item}{5}{B}
The height of a launched weather balloon can be represented by an exponential model, (y=a\cdot b^x), where (y) is the height in km after (x) min. After 2 min, the height is 3 km, and after 3 min it is 9 km. Which represents another point on the graph of this model?
A. ((1, -3))
B. ((1, 1))
C. ((4, 15))
D. ((4, 16))
\answer{5B. B}
\end{ta-item}
\begin{ta-item}{6}{B}
A school of fish has a population of 200. The population is decreasing at a rate of 5\% per year. Select the model of the population after (t) years that reveals the quarterly decay factor.
A. (y=200(0.95)^t)
B. (y=200(0.9875)^{4t})
C. (y=200(0.9875)^t)
D. (y=200(0.9873)^{4t})
\answer{6B. D}
\end{ta-item}
\begin{ta-item}{7}{B}
What function is the inverse of the exponential function (y=\left(\frac{1}{3}\right)^x)?
A. (y=x^{\frac{1}{3}})
B. (y=\left(\frac{1}{3}\right)^x)
C. (y=\log_{\frac{1}{3}} x)
D. (y=\log_x\left(\frac{1}{3}\right))
\answer{7B. C}
\end{ta-item}
\begin{ta-item}{8}{B}
What is the value of (n) in the formula (S=\frac{a_1(1-r^n)}{(1-r)}) for the sum of the geometric series (0.4+1.2+3.6+\ldots+7873.2)?
(n=) [ ]
\answer{8B. 10}
\end{ta-item}
\begin{ta-item}{9}{B}
What is the solution to the equation (\log_3(2x-3)=-1)? Round to the nearest thousandth.
(x=) [ ]
\answer{9B. 1.667}
\end{ta-item}
\begin{ta-item}{10}{B}
The graph shows the function (f(x)=32\left(\frac{1}{2}\right)^x-2). What is the value of the inverse function (f^{-1}) at (x=2)?
\placeholder{graph}{Prompt graph (not part of the answer), to the right of the text: a decreasing exponential curve on a gray grid. Axes labeled x and y with arrowheads, origin O. x-axis labels 4 and 6, gridlines every 1 unit, window about x = -2 to 7. y-axis labels 10, 20, 30, gridlines every 5 units, window about y = -10 to 35. Labeled points on the curve: (0, 30), (1, 14), (2, 6), (3, 2), (4, 0). The curve leaves the top of the grid at the y-axis (arrowhead) and falls through (4, 0) to below the x-axis, flattening toward its asymptote y = -2 and ending in an arrowhead to the right near x = 6.}
(f^{-1}(2)=) [ ]
\answer{10B. 3}
\te{The Form B docx reads "inverse function, f −1 at x = 2" with a comma; the TE page has no comma (Form A has "inverse function, f^{-1}, at x = 2"). TE image followed.}
\end{ta-item}
\begin{ta-item}{11}{B}
What is the asymptote of the graph of (g(x)=\log_2(x+4))?
A. vertical at (x=0)
B. vertical at (x=-4)
C. horizontal at (y=0)
D. horizontal at (y=-4)
\answer{11B. B}
\end{ta-item}
\begin{ta-item}{12}{B}
Complete the equation of the inverse of the function (f(x)=\log_4(3x)).
\placeholder{figure}{Tile bank of four drag tiles, left to right: 1/3, 3, 4, 1/4.}
(f^{-1}(x)=) [ ]([ ])^x
\answer{12B. (f^{-1}(x)=\frac{1}{3}(4)^x): first box 1/3, second box 4}
\te{The earlier Topic6AssessFormB.tex orders the tiles 4, 1/3, 1/4, 3; the TE page orders them 1/3, 3, 4, 1/4. TE image followed.}
\end{ta-item}
\begin{ta-item}{13}{B}
Use the properties of logarithms to choose the expression (x\log 2-\log\frac{1}{3}) as a single logarithm.
A. (\log\left(2x-\frac{1}{3}\right))
B. (\log\left(\frac{2}{3x}\right))
C. (\log\left(2^x-\frac{1}{3}\right))
D. (\log(3(2)^x))
\answer{13B. D}
\end{ta-item}
\begin{ta-item}{14}{B}
The profit for a new website, in thousands of dollars, can be represented by (5^x), where (x) is time in months. Find the time when the company's profit is \$15,000. Select the exact solution as a logarithm, and select an approximate solution.
A. (x=\log_{15}5)
B. (x=\log_5 15)
C. (x=\log 15)
D. (x=1.689)
E. (x=0.594)
F. (x=0.176)
\answer{14B. B, D}
\te{computed: log_5 15 = ln 15 / ln 5 = 1.6826..., so the exact choice B rounds to 1.683, not the printed 1.689 of D. The printed key (B, D) stands; 1.689 is printed the same in the Realize PDF and the docx. Distractors: E = log_15 5 = 0.594 (correct); F = 0.176 = log 1.5, while log 15 = 1.176 (Form A's parallel F, 0.903, equals log 8).}
\te{Choices print in two columns: A, B, C on the left and D, E, F on the right.}
\end{ta-item}
\begin{ta-item}{15}{B}
What is the sum of the geometric series (S=1.375+1.1+0.88+\ldots+0.031)? Round to the nearest thousandth.
(S=) [ ]
\answer{15B. 6.751}
\te{computed: r = 0.8; 0.031 is the rounded 18th term (1.375(0.8)^{17} = 0.03096); S_{18} = 6.875(1-0.8^{18}) = 6.7512, agrees with the key.}
\end{ta-item}
\begin{ta-item}{16}{B}
Complete the steps to write the expression (2\log a+2\log b-3\log c) as a single logarithm using the properties of logarithms.
\placeholder{figure}{Tile bank of three drag tiles, left to right: Product, Quotient, Power.}
(2\log a+2\log b-3\log c)
(=\log a^2+\log b^2-\log c^3) [ ]
(=\log a^2b^2-\log c^3) [ ]
(=\log\left(\frac{a^2b^2}{c^3}\right)) [ ]
\answer{16B. Power; Product; Quotient (boxes top to bottom)}
\end{ta-item}
\begin{ta-item}{17}{B}
For the function (f(x)=2\cdot(3)^{-x}), identify the (y)-intercept and asymptote.
(y)-Intercept: [ ]
Asymptote: (y=) [ ]
\answer{17B. y-Intercept: 2; Asymptote: y = 0}
\end{ta-item}
\te{The earlier Lesson_planning/Topic6AssessFormB.tex is not the full Form B: it keeps 8 of the 17 items, renumbered 1-8 (Form B items 9, 10, 7, 11, 12, 13, 14, 16 in that order), with a key that agrees with the TE marks for those items.}
\end{document}
